\documentclass[11pt,a4paper]{article}
\usepackage[margin=25mm]{geometry}
\usepackage[T1]{fontenc}
\usepackage{lmodern}
\usepackage{microtype}
\usepackage{amsmath}
\usepackage{booktabs,tabularx,array}
\usepackage{graphicx}
\usepackage{xcolor}
\usepackage[hidelinks]{hyperref}
\setlength{\parindent}{0pt}
\setlength{\parskip}{6pt}
\newcolumntype{Y}{>{\raggedright\arraybackslash}X}

% Individual submission. Author details confirmed by the student.
\newcommand{\studentname}{DAM MINH TIEN}
\newcommand{\studentid}{A0355091E}

% Draft instructions are intentionally visible. Replace each box with measured
% results or final prose; no draft instruction should remain in the submitted PDF.
\newcommand{\draftnote}[1]{\par\begingroup\small\color{blue!55!black}
  \noindent\textbf{Draft instruction:} #1\par\endgroup}
\newcommand{\figureplaceholder}[2]{%
  \begin{figure}[htbp]\centering
  \fbox{\parbox[c][30mm][c]{0.9\linewidth}{\centering\small #1}}
  \caption{#2}\end{figure}}

% EDA figures appear only when the actual exported artifacts are present.
\newcommand{\edaplot}[2]{%
  \begin{figure}[htbp]\centering
  \IfFileExists{figures/#1.pdf}{%
    \includegraphics[width=0.94\linewidth]{figures/#1.pdf}%
  }{%
    \fbox{\parbox[c][20mm][c]{0.9\linewidth}{\centering\small
      Awaiting measured EDA output: \texttt{\detokenize{#1}}}}%
  }%
  \caption{#2}\label{fig:#1}\end{figure}}

\begin{document}
\begin{titlepage}
  \centering
  \vspace*{20mm}
  {\Large National University of Singapore\par}
  \vspace{8mm}
  {\large DSA5208\par}
  \vspace{18mm}
  {\LARGE\bfseries Project 2\par}
  \vspace{5mm}
  {\LARGE Singapore Rainfall Analysis and Prediction\par}
  \vspace{20mm}
  {\large \studentname\par}
  \vspace{3mm}
  {\large Student ID: \studentid\par}
  \vspace{5mm}
  {\large Individual submission\par}
  \vfill
  {\large 18 October 2026\par}
\end{titlepage}

\pagenumbering{roman}
\section*{Abstract}
\draftnote{Write 150--200 words after completing the results. State the two
assignment tasks, the 2017--2024 rainfall data, the actual Spark/MLlib approach,
one quantified Task 1 finding, the selected model's final-test performance, and
one material limitation. Do not report planned experiments as completed work.}
\tableofcontents
\clearpage
\pagenumbering{arabic}

\section{Introduction}
\label{sec:introduction}
This project analyses rainfall measurements from Singapore monitoring stations
for 2017--2024~\cite{dataset}. Task 1 examines rainfall patterns within a year and across these
years. Task 2 estimates the probability that a station will accumulate more
than a specified rainfall threshold in the next 30 minutes. The assignment
requires structured data processing with Apache Spark and model development
and evaluation with Spark MLlib.
\draftnote{Add a short paragraph describing the actual completed approach and
scope. Keep this section to about half a page. Separate assignment requirements
from implementation choices; no separate research questions or novelty claims
are needed.}

\section{Dataset and Preprocessing}
\label{sec:data}
\subsection{Data source and record structure}
\draftnote{Identify the NEA annual datasets for 2017--2024, retrieval date,
files, units, stations and Singapore timezone. Describe the fields used.
Explain observation time, update time and the verified meaning of a five-minute
reading. State any unresolved interpretation as an assumption. Cite the dataset
entry and any documentation actually used.}

\subsection{Cleaning and quality assessment}
\draftnote{Explain exact duplicates, conflicting revisions, invalid values,
missing readings, gaps and station metadata. Missing rainfall must remain
distinct from zero rainfall. State the cleaning order and exclusion counts so
that overlapping problems are not counted twice.}
\begin{table}[htbp]\centering\small
\caption{Annual data quality after applying the stated cleaning rules.
Replace draft cells with actual counts and define the coverage denominator.}
\label{tab:quality}
\begin{tabularx}{\linewidth}{lYYYYY}\toprule
Year & Raw rows & Retained rows & Stations & Coverage & Excluded rows\\\midrule
2017--2024 & \multicolumn{5}{l}{[Insert one row per year from the Spark audit]}\\
\bottomrule\end{tabularx}
\end{table}

\subsection{Analysis cohorts and processed outputs}
\draftnote{Define valid station-days and the comparable station cohort for
annual analysis. Report attrition from cleaned observations to usable history,
complete future windows and prediction samples. Explain the CSV-to-Parquet
workflow and which cleaned outputs are reused by both tasks.}

\section{Computing Environment}
\label{sec:environment}
Project code is executed from Jupyter notebooks, which retain the configuration,
outputs and Spark application/job identifiers. Reusable Python functions are
imported into the notebook cells. The local environment preflight has run;
distributed processing remains pending the billing linkage described in
Appendix~\ref{app:preflight}.
\draftnote{Describe the environment actually used. Summarise Spark processing,
worker execution, storage and MLlib training. State measured processing/training
times and relevant constraints. Use about half a page; put long configs/logs in
the appendix. Cloud resources remain conditional on the approved billing and
credit checks; a proposed cluster is not an observed run. For every submitted
notebook action, review its Spark Web UI after the generated job(s)
reach a final state. Use the live application UI if a notebook/application is
still open, or Spark History Server if the application has ended. Review the
Jobs summary and the detail page for each Spark Job ID. Capture
and retain screenshots showing the matching application/job IDs, final status,
duration and available stages/tasks; list and embed every screenshot in the
appendix. Confirm event logging is enabled before runs so completed applications
can be reopened. Do not use the Dataproc Jobs page as a substitute for Spark UI
evidence~\cite{sparkui,dataprocui}.}
\begin{table}[htbp]\centering\small
\caption{Actual computing environment for the reported run.}
\label{tab:environment}
\begin{tabularx}{\linewidth}{lY}\toprule
Item & Actual value\\\midrule
Platform and region & [To be recorded]\\
Master and workers & [VM types, number, vCPU and RAM]\\
Storage & [Disk configuration and input/output storage]\\
Runtime & [Dataproc image; Spark, Python and Java versions]\\
Execution settings & [Relevant executor/partition settings]\\
Run identifier & [Run ID and code commit]\\
\bottomrule\end{tabularx}\end{table}

\section{Task 1: Rainfall Patterns}
\label{sec:task1}
\subsection{Scope and definitions}
The exploratory data analysis (EDA) addresses the assignment's two descriptive
requirements: rainfall patterns within a year and variation across 2017--2024.
The implementation first analyses 2017--2023. The full eight-year analysis is
reserved for execution after the Task 2 configuration and model are frozen.
This order prevents findings from the test year from influencing feature or
model selection. The figures and numerical findings in this section remain
pending execution of the distributed EDA notebook.

A valid reading is a retained, unique station--timestamp measurement on the
exact five-minute grid. A \emph{wet} reading has rainfall strictly greater than
0 mm; a valid zero denotes a dry interval. Missing and rejected records are
excluded from rainfall statistics and retained in the coverage denominator.
Following the accepted project assumption, timestamps identify the ends of
five-minute intervals; this convention has not been confirmed by NEA. Month,
date and hour groups refer to the interval-end timestamps in Singapore time
(UTC+8). Consequently, calendar-labelled totals describe readings ending in
that period, including the interval ending at midnight, rather than an
independently verified meteorological daily-total convention.

\subsection{Coverage and comparable station cohorts}
For station $s$ and a calendar period $P$, let $N_{s,P}$ be the number of valid
readings and $E_P=288D_P$ the expected count for $D_P$ calendar days. Coverage is
\begin{equation}
 C_{s,P}=\frac{N_{s,P}}{288D_P}.
\end{equation}
Leap days are included. The fixed reference network is the set of station IDs
in the cleaned archive. Its zero-observation station-periods remain visible;
commissioning and retirement dates are unavailable, so low fixed-calendar
coverage cannot be interpreted directly as instrument downtime. The coverage
heatmap, excluded-row audit and gap summaries provide the context for the
rainfall comparisons.

The primary monthly summaries require station-month coverage of at least
95\%. Hourly summaries additionally require at least 95\% of the expected
readings for that hour within the station-month: $12D_P$ readings, since each
hour contains twelve five-minute intervals. Group counts and valid-observation
counts accompany each summary. These checks reduce comparisons based on sparse
records but do not establish that missing observations are unrelated to rain.

Annual comparisons use the same stations throughout the analysed years. The
implementation reports the number of common stations at fixed coverage cutoffs
of 90\%, 95\% and 99\%, including a count of zero when a cohort is empty. A
cutoff is not relaxed in response to a test-year finding. An all-observed-station
view is exported separately to reveal how changing network composition affects
the descriptive totals; it cannot replace a common-cohort comparison without
stating that limitation.
\edaplot{coverage_heatmap}{Valid five-minute readings divided by full-calendar
expected readings, by station and year. Zero-observation years remain in the
reference network. Coverage describes data availability, not verified station
operating time.}

\subsection{Rainfall distribution and large readings}
The distribution summary separates the frequency of rain from its intensity.
For a stated group of valid readings $\mathcal V$, the wet frequency and
conditional wet intensity are
\begin{equation}
 f_{\mathrm{wet}}=
 \frac{\sum_{i\in\mathcal V}\mathbf{1}(r_i>0)}{|\mathcal V|},
 \qquad
 \mu_{\mathrm{wet}}=
 \frac{\sum_{i\in\mathcal V}r_i\mathbf{1}(r_i>0)}
 {\sum_{i\in\mathcal V}\mathbf{1}(r_i>0)}.
\end{equation}
Conditional intensity is undefined for a group with no wet readings. The
notebook exports valid, dry and wet counts; mean rainfall per valid interval;
wet-reading median, P90, P95 and P99; and the maximum observed reading for each
year. Quantiles are approximate Spark aggregates with accuracy parameter
10,000. The annual distribution summaries pool valid readings and therefore
give greater weight to stations with more observations. This weighting differs
from the equal-station summaries used for the primary seasonal comparisons.

A fixed-bin histogram shows wet rainfall in mm per five-minute interval;
a logarithmic vertical scale makes the sparse upper tail visible. Dry records
are shown as a separate proportion and excluded from the wet histogram's
normalisation. The reference for unusually large readings is the wet-reading
P99 computed using 2017--2022 only. Exceedance frequencies use this unchanged
reference in every year. It is a descriptive threshold, not an official warning
level. Large valid readings are retained; the maximum alone cannot establish
that an event is an error or a meteorological extreme.
\edaplot{rainfall_distribution}{Annual wet-record percentages and the
conditional distribution of positive five-minute rainfall. Denominators are
valid readings and wet readings, respectively; both panels pool observations.
The wet-bin proportions are displayed on a logarithmic scale.}

\subsection{Patterns within a year}
For an eligible station-month, observed rainfall is
$T_{s,P}=\sum_{i\in\mathcal V_{s,P}}r_i$ in mm. Monthly rainfall is the arithmetic
mean of these station totals; wet frequency and conditional wet intensity are
also averaged across eligible stations. This averaging assigns each station
one weight within a month. The eligible station set may vary between months,
so the support panel must be read alongside a seasonal peak or trough.
Observed totals are not divided by coverage to reconstruct missing rainfall,
and rainfall from different stations is not summed into a Singapore-wide total.
Even a qualifying month with 95\% coverage may omit part of its rainfall.

The monthly analysis compares amount, frequency and conditional intensity.
These quantities distinguish a month with more wet intervals from one with
larger readings when rain occurs. A complementary daily analysis retains only
station-dates with all 288 unique grid readings, computes their observed totals,
and averages daily rainfall within each station-month before averaging across
stations. The corresponding support table reports complete station-days relative
to both observed and full-calendar station-days. Complete-day filtering can
exclude periods with missing observations during rain; it is a sensitivity
view, not a correction for missingness.

The month--hour heatmap averages wet frequencies equally across eligible
station--year--month--hour groups. Separate yearly hourly curves average eligible
station-month groups within a year. Thus neither plot allows a station with more
valid observations to dominate by its record count. Support is recorded by hour,
and the time axis labels the end of the five-minute interval in Singapore time.
\edaplot{monthly_hourly_patterns}{Mean observed station-month rainfall (mm)
and equal-group wet frequency by month and interval-end hour (SGT). Monthly
coverage and, for hourly comparisons, hour coverage must each be at least 95\%.
Missing observations contribute no rainfall value. Eligible station/group
counts are exported with the plotted statistics.}
\draftnote{Use the generated monthly and hourly extrema with their support
counts to state the highest and lowest represented periods, the rainfall
difference in mm or wet-frequency difference in percentage points, and whether
the pattern recurs across years. A missing month/hour is unavailable evidence,
not a zero. Discuss any changing station support before interpreting the peak.
Place the frequency/intensity, complete-day and yearly-hourly panels in the
appendix when they clarify the main pattern.}

\subsection{Changes across 2017--2024}
Annual observed rainfall and wet frequency are compared within the common
station cohort. Each station receives equal weight. Coverage sensitivity
reports the station count and the highest--lowest year difference in observed
rainfall at all three fixed cutoffs. Consistent comparisons across these cohorts
provide stronger descriptive evidence than an unadjusted network-wide sum.
If a cutoff yields no common stations, that comparison is unavailable. The
changing-cohort view remains descriptive and is labelled accordingly. No trend
line or causal attribution is inferred solely from these eight annual values.
\edaplot{annual_patterns}{Equal-station mean observed annual rainfall (mm)
within stations meeting the stated coverage cutoff in every analysed year.
Curves use 90\%, 95\% and 99\% cutoffs, with common-station counts. Empty cohorts
are reported explicitly. Totals remain observed values, not imputed full-year
rainfall.}
\draftnote{Insert the measured annual differences and common-station counts.
Compare year rankings across cutoffs and with the all-observed-station view.
Explain any disagreement through the measured coverage/cohort changes; do not
claim a long-term climate trend.}

\subsection{Station differences and observed wet spells}
The station-coordinate plots describe wet frequency and exceedance of the fixed
training P99. Only stations whose recorded coordinates remain unchanged across
the analysed period are mapped. Annual summaries require the primary coverage
cutoff and are averaged equally over each station's eligible years. The number
of eligible years and valid readings accompanies each point; different stations
may represent different years. The plots show observed station measurements
without spatial interpolation, so they do not estimate rainfall over an
unmonitored area or establish that location causes a difference.

An observed wet spell is a maximal sequence of positive readings at exactly
300-second spacing for one station. Its observed duration is five minutes times
the number of readings. A gap terminates the observed sequence; it does not
prove that rain stopped. Spells adjoining a gap or an analysis boundary are
flagged as censored and summarised separately from spells bounded by valid dry
readings. The notebook reports counts, duration median/P90/P99 and maximum
observed duration by the year of spell start. These are station-level observed
sequences, not identified meteorological storms. Spatial, tail and spell figures
are supplementary evidence and belong in the appendix when discussed.

\subsection{Summary of Task 1 findings}
\draftnote{After the Spark run, replace this note with two or three quantified
findings supported by the exported CSV tables and figures: one within-year
pattern, one comparison across years, and a material coverage or distribution
result. Give units, cohort/support and the relevant limitation. Connect measured
time patterns to time features and observed short-term persistence to rainfall
history features in Task 2; retain the pre-test feature design. Every numerical
statement must match the notebook's computed outputs.}

\section{Task 2: Rainfall Probability Prediction}
\label{sec:task2}
\subsection{Prediction target and input contract}
For station $s$, prediction time $t$ and rainfall threshold $\tau$ in mm, the
requested output is
\begin{equation}
  p(s,t,\tau)=\Pr\!\left(R_{30}(s,t)>\tau\mid\mathcal{H}_{s,t}\right),
\end{equation}
where $\mathcal{H}_{s,t}$ denotes the information available at prediction time.
\draftnote{Once interval semantics are verified, define $R_{30}$ explicitly as
the sum of the six aligned future five-minute intervals. Explain the strict
greater-than comparison, interval boundaries, label completeness and feature
availability. State threshold domain and behaviour for gaps, unseen stations
and unsupported queries. Identify retrospective availability assumptions.}

\subsection{Feature selection}
\draftnote{Describe only features actually used. Explain why each group is
useful, how history is calculated and how future observations/revisions are
excluded. Connect time/history features to Task 1 findings when supported.}
\begin{table}[htbp]\centering\small
\caption{Feature groups in the final model. Replace candidate descriptions
with the exact implemented features and windows.}
\label{tab:features}
\begin{tabularx}{\linewidth}{lYY}\toprule
Group & Features and window & Reason and availability\\\midrule
Station and time & [Encoding; calendar fields] & [Location/season/time patterns]\\
Threshold & [Requested rainfall threshold] & [Defines the exceedance event]\\
Rainfall history & [Actual past windows] & [Recent rainfall information]\\
Missingness & [Actual validity indicators] & [Incomplete observation history]\\
\bottomrule\end{tabularx}\end{table}

\subsection{Temporal split and training procedure}
\draftnote{The table below is the planned split; update it if the documented
protocol changes before final testing. Report actual sample counts and
prevalence. Explain the 30-minute boundary purge and train-only preprocessing.
Keep model selection, calibration fitting and final-test evaluation separate.}
\begin{table}[htbp]\centering\small
\caption{Temporal evaluation split. Counts are measured after cleaning,
history/label eligibility checks and boundary purging.}
\label{tab:split}
\begin{tabularx}{\linewidth}{lYYY}\toprule
Subset & Period & Purpose & Actual support\\\midrule
Training & 2017--2022 & Fit preprocessing and models & [Base/expanded counts]\\
Selection & Jan--Jun 2023 & Select model/settings & [Counts and prevalence]\\
Calibration & Jul--Dec 2023 & Fit calibration mapping & [Counts and prevalence]\\
Final test & 2024 & Frozen final evaluation & [Counts and prevalence]\\
\bottomrule\end{tabularx}\end{table}

\subsection{Baseline and Spark MLlib models}
\draftnote{Describe the empirical climatology baseline, Logistic Regression
and Gradient-Boosted Trees. State the actual preprocessing, limited tuning
search, selected settings and validation criterion. Explain probability output
and calibration when applied. Full hyperparameters belong in the appendix.
Do not select the final model by its 2024 score.}

\subsection{Predictive performance}
\draftnote{Compare models on the same final-test cohort and frozen threshold
grid. Define PR-AUC, ROC-AUC and Brier briefly and explain their roles. Distinguish
Spark areaUnderPR from Average Precision. Define macro aggregation and metric
settings. Report support and N/A where a group contains only one class.}
\begin{table}[htbp]\centering\small
\caption{Final-test model comparison on the same cohort. Specify the
threshold aggregation, support and calibration status in the final caption.}
\label{tab:models}
\begin{tabularx}{\linewidth}{Yccc}\toprule
Model & PR-AUC & ROC-AUC & Brier\\\midrule
Climatology & [Result] & [Result] & [Result]\\
Logistic Regression & [Result] & [Result] & [Result]\\
Gradient-Boosted Trees & [Result] & [Result] & [Result]\\
\bottomrule\end{tabularx}\end{table}

\begin{table}[htbp]\centering\small
\caption{Per-threshold final-test performance. Include baseline comparison
and the selected model's metrics, with the same support for both.}
\label{tab:thresholds}
\begin{tabularx}{\linewidth}{lYYYYY}\toprule
$\tau$ (mm) & Samples & Positives & Prevalence & PR-AUC & Brier\\\midrule
Grid: [values] & [Count] & [Count] & [Fraction] & [Result] & [Result]\\
\bottomrule\end{tabularx}\end{table}
\figureplaceholder{PR curves at representative thresholds fixed before testing.}
{[Draft] Precision--recall curves on the final-test cohort. Identify thresholds,
models, positive prevalence, support and evaluator settings.}
\draftnote{Explain the comparison that matters: which validation-selected
model performs well on the test, where it exceeds or falls below climatology,
and how rare events change interpretation. Do not restate every table cell.}

\subsection{Calibration and probability consistency}
\figureplaceholder{Reliability diagram with counts per bin; raw/calibrated if fitted.}
{[Draft] Predicted probabilities and observed exceedance frequencies on the
final test. State threshold, binning, support and calibration fit period.}
\draftnote{Discuss whether probabilities match observed frequencies and whether
calibration improved measured probability quality. Report monotonicity across
thresholds for the same station/time/history. State limitations for unseen
thresholds and any correction actually included in inference.}

\subsection{Feature contribution and error analysis}
\draftnote{Summarise the A/B comparison: station/time/threshold versus adding
history/missingness, using the same samples. Explain a few supported failure
patterns such as onset versus continuation, heavy-rain thresholds or sparse
stations. Give support and uncertainty where available. Keep extensive strata
tables in the appendix.}

\section{Discussion and Limitations}
\label{sec:discussion}
\draftnote{Connect Task 1 patterns with Task 2 behaviour. Discuss only material
limits demonstrated or applicable to this run: rainfall-only information,
missingness/revisions, availability assumptions, rare events, eight-year scope,
unseen stations and supported threshold range. Explain how each affects the
conclusions. Report disappointing results accurately; additional work is optional.}

\section{Reproducibility and Conclusion}
\label{sec:conclusion}
\subsection{Reproducing the submitted work}
\draftnote{Identify the submitted code, README, manifest, exact runtime/config,
model package and entry point. State the actual fresh-session reload/prediction
check and its result once performed. The package should not depend on access to
a personal cloud bucket.}
\subsection{Conclusion}
\draftnote{Summarise one or two quantified findings from each task and the
conditions under which the selected model's probabilities are useful. Introduce
no new experiment, metric or claim.}

\section*{AI Usage Declaration}
\addcontentsline{toc}{section}{AI Usage Declaration}
\draftnote{Use first-person singular for this individual project. Complete the
course's declaration with the actual tool/model names and uses across planning,
coding, debugging, analysis and report preparation. Revise it as the work evolves.}
I used [AI tool and model] to [specific uses]. I am responsible for the content
and quality of the submitted work.

\begin{thebibliography}{9}
\bibitem{dataset}
National Environment Agency. \emph{Historical Rainfall across Singapore}.
\url{https://data.gov.sg/collections/2279/view}.
\bibitem{sparkui}
Apache Spark. \emph{Web UI}.
\url{https://spark.apache.org/docs/latest/web-ui.html}.
\bibitem{dataprocui}
Google Cloud. \emph{Managed Service for Apache Spark Component Gateway}.
\url{https://docs.cloud.google.com/managed-spark/docs/concepts/accessing/gateways}.
\end{thebibliography}
\draftnote{Add verified references for the Spark version, methods and any
meteorological interpretation actually cited. Record dataset retrieval date.
Do not include sources that were not used.}

\appendix
\section{Additional Results and Run Details}
\draftnote{Include only useful supplementary material: complete threshold
tables, coverage sensitivity, full model settings, selected runtime evidence
and representative inference queries. Main text must remain understandable
without this appendix. Omit empty appendix subsections from the final report.}

\subsection{Notebook execution and billing preflight}
\label{app:preflight}
The sole authorised credit is \texttt{3LU63BPGY77JFDQQ}, with a user-provided
original and remaining value of USD~50. The authorised billing account is
Billing Account for Education, ID \texttt{01DAD8-AE5F1A-E3910A}.
The credit ID and billing account ID identify separate objects. Their live
association and remaining credit must be checked before paid execution.
The local \path{notebooks/00_environment.ipynb} preflight on 11 October 2026
found project \texttt{new-cloud-256902} billing-disabled. Linking it to Education
was denied because the active identity lacked project billing-assignment
permissions. No cloud bucket, cluster or Spark job was created by this run.
\draftnote{Replace this preflight status with the actual resolved environment
and dated billing evidence when available. Retain executed notebooks with
visible, sanitised outputs. Run the audit through
\path{notebooks/01_data_quality.ipynb}, and record the notebook/cell, input
manifest hash, configuration, application ID and every Spark Job ID.}

\subsection{Spark Web UI evidence for completed jobs}
\label{app:spark-ui-evidence}
\draftnote{For each notebook action, open its completed
application in the live Spark Web UI if it remains open, or in Spark History
Server after it ends, and review the Jobs tab. The summary shows job status and
duration; each job detail page provides its stages, tasks and DAG as
available~\cite{sparkui}. Include the original screenshot of the application
Jobs summary and a detail screenshot for every Spark Job ID. Preserve visible
Application ID, Job ID and final status; do not use Dataproc Jobs or YARN pages
as substitutes. Store screenshots under
\path{report/evidence/spark-ui/<RUN_ID>/} and maintain
\path{report/evidence/spark-ui/index.csv}, with one record per submitted Spark
notebook action, its application and Spark Job ID(s), final status, duration and all
screenshot paths. Record failed or cancelled jobs accurately. If an application
created no Spark Job ID, retain its summary screenshot and state this explicitly.
If no job has run, do not invent screenshots or results.}
\begin{table}[htbp]\centering\scriptsize
\caption{Spark Web UI evidence for each submitted notebook action.}
\label{tab:spark-ui-evidence}
\begin{tabularx}{\linewidth}{p{1.7cm}YYYYY}\toprule
Notebook action/cell & Application ID & Spark Job ID(s) & Final status & Duration & Screenshot paths\\\midrule
Row: [one per notebook action] & [Application ID] & [All Job IDs] & [Actual state(s)] & [Actual time] & [Summary PNG; every job-detail PNG]\\
\bottomrule\end{tabularx}
\end{table}
\draftnote{Embed each indexed original screenshot in this appendix and keep
text readable at page size. The index must account for every submitted Spark
notebook action and all of its Spark Job IDs; screenshots are evidence of review, not
substitutes for logs, outputs or measured results.}
\end{document}
